\scanpage{115}\mspage{\handwritten{12}}

\inserted{\handwritten{implies}} inductively that an automorphism of the system of $\underline{F}_i$'s related by the source and target maps $\underline{s}_{\unsure{1|l}}^i,\underline{t}_1^i$ is the identity. Thus it is correct afte\unsure{r} all, it seems, that the category structure of a coherator implies all other structure relevant to us.
% slanted
\margin{\handwritten{Maybe false, \unsure{we} should give moreover the subcategory $\underline{C}_{\unsure{o|0}}$ perhaps}}

I do believe that the description given so far of what I mean by a coherator, \deleted{\unsure{acting a}} namely something acting like a kind of pattern in order to define a corresponding notion of ``stacks'' (which in turn should be the basic coefficient objects in non commutative homological algebra, as well as a convenient description of homotopy types) embodies some of the essential features of the theory\inserted{ still in embryo} that wants to be developped. It is quite possible of course that some features are lacking still, for instance that some extra conditions have to be imposed upon $\underline{C}$, possibly of a very different nature from mere irredundancy conditions (which, I feel, are kind of irrelevant in this set-up). Only by pushing ahead and working out at least in outline the main aspects of the formalism of stacks, will it become clear whether or not extra conditions on $\underline{C}$ are needed. I would like at least to make a commented list of these main aspects, and possibly do some heuristic pondering on some of these in the stride, or afterwards. For today it seems a little late though -- I have been pretty busy with non mathematical work most part of the day, and the next two days I'll be busy at the university. Thus I guess I'll send off this letter tomorrow, and send you \unsure{an}\inserted{ later} elaboration (presumably much in the style of this unending letter) if you are interested. In any case I would appreciate any comments you make -- that's why I have been writing you after all ! I will probably send copies to Ronnie Brown, Luc Illusie and Jean Giraud, in case they should be interested (I guess at least Ronnie Brown is). Maybe the theory is going to take off after all, in the long last !

Very cordially yours\\
\hfill\handwritten{Alexander}\\

PS (23.2) I noticed a rather silly mistake in the notes of two days ago,when stating that the categories $\underline{C}_n$ admit fiber\deleted{ed} products -- what is true is that the category $\underline{C}_o$ has fiber\deleted{ed} products (by construction\inserted{ practically}), and that these are fiber\deleted{ed} products also in the categories $\underline{C}_n$ (by construction equally) i.e. that the inclusion functors $\underline{C}_o\handwritten{\to}\underline{C}_n$ commute to fiber\deleted{ed} products. Stacks in a category $C$ correspond to functors $\underline{C}_{\infty}\handwritten{\to}C$ whose restriction to $\underline{C}_o$ commutes with fiber products. I carried the mistake along in the \deleted{notes} yesterday notes -- it does'nt really change anything substantially. I will have to come back anyhow upon the basic notion of \unsure{a coherator}

% archive sheet number: 114
